\section*{Hoofdstuk \ref{ch:b2:chemical-potential} --- Chemische potentiaal en mengsels}

\begin{solution}{exo:b2:chemical-potential:1}
De raaklijn bij $x_2 = 0.4$ snijdt de randen in $\bar V_1 \approx 14$ en $\bar
V_2 \approx 49$ (modeleenheden). Dan is $0.6 \times 14 + 0.4 \times 49 = 28.0$,
het molair volume dat op de kromme bij $x_2 = 0.4$ wordt afgelezen, onder de
$0.6 \times 18 + 0.4 \times 58 = 34.0$ van de zuivere volumes.
\end{solution}

\begin{solution}{exo:b2:chemical-potential:2}
$n(\text{water}) = 1000/18.015 = \qty{55.5}{mol}$, $x_1 = 55.5/56.5 = 0.982$;
$p = 0.982 \times 3.17 = \qty{3.11}{kPa}$: een verlaging met 1.8\,\%.
\end{solution}

\begin{solution}{exo:b2:chemical-potential:3}
$p(\ce{O2}) = 0.2095 \times \qty{1.013e5}{Pa} = \qty{2.12e4}{Pa}$;
$c = 1.3 \times 10^{-5} \times 2.12 \times 10^4 = \qty{0.28}{mol/m^3} =
\qty{2.8e-4}{mol/L}$, dat is $2.8 \times 10^{-4} \times 32.0 \times 10^3 =
\qty{8.8}{mg/L}$.
\end{solution}

\begin{solution}{exo:b2:chemical-potential:4}
$Q = (1.0)^2/[(1.0)(0.010)^3] = 1.0 \times 10^6$; $RT\ln Q = 2.479 \times 13.82 =
\qty{34.2}{kJ/mol}$; $\Delta_r G = -32.8 + 34.2 = \qty{+1.4}{kJ/mol} > 0$: in dit
mengsel, zo arm aan waterstof, ontleedt ammoniak.
\end{solution}

\begin{solution}{exo:b2:chemical-potential:5}
Gibbs--Duhem: $x_1\dd\bar V_1 + x_2\dd\bar V_2 = 0$, dus $\dd\bar V_2 =
-(0.70/0.30) \times (-0.30) = \qty{+0.70}{mL/mol}$.
\end{solution}

\begin{solution}{exo:b2:chemical-potential:6}
$c(\text{deeltjes}) = 2 \times 9.0/58.44 = \qty{0.308}{mol/L} =
\qty{308}{mol/m^3}$; $\Pi = 308 \times 8.314 \times 310 = \qty{7.9e5}{Pa}$,
ongeveer \qty{7.9}{bar}. Het inwendige van een rode bloedcel heeft dezelfde \omterm{def:b2:chemical-potential:osmotic-pressure}{osmotische druk}:
geen netto waterstroom. In zuiver water dringt water de cel binnen, die opzwelt en
barst.
\end{solution}

\begin{solution}{exo:b2:chemical-potential:7}
$c = \Pi/RT = 825/(8.314 \times 298.15) = \qty{0.333}{mol/m^3}$; in
\qty{1.000e-4}{m^3} is $n = \qty{3.33e-5}{mol}$, $M = 2.00/3.33 \times 10^{-5}
= \qty{6.0e4}{g/mol}$ (\qty{60}{kg/mol}). $\Delta T_f = 1.86 \times 3.33
\times 10^{-5}/0.100 = \qty{6.2e-4}{K}$: onmeetbaar, terwijl de \omterm{def:b2:chemical-potential:osmotic-pressure}{osmotische
druk} overeenkomt met een waterkolom van \qty{8}{cm}.
\end{solution}

\begin{solution}{exo:b2:chemical-potential:8}
(a) $I = \qty{1.0e-3}{mol/kg}$, $\log\gamma_\pm = -0.51 \times 0.0316$,
$\gamma_\pm = 0.964$. (b) $I = \frac12(4 \times 1.0 + 1 \times 2.0) \times
10^{-3} = \qty{3.0e-3}{mol/kg}$, $\log\gamma_\pm = -0.51 \times 2 \times 0.0548$,
$\gamma_\pm = 0.879$. (c) $I = \frac12(4 + 4) \times 10^{-3} = \qty{4.0e-3}{mol/kg}$,
$\log\gamma_\pm = -0.51 \times 4 \times 0.0632$, $\gamma_\pm = 0.74$.
\end{solution}

\begin{solution}{exo:b2:chemical-potential:9}
$b = 0.310/1.86 = \qty{0.167}{mol/kg}$; $n = 0.167 \times 0.0500 =
\qty{8.33e-3}{mol}$; $M = 1.50/8.33 \times 10^{-3} = \qty{180}{g/mol}$
(bijvoorbeeld een hexose).
\end{solution}

\begin{solution}{exo:b2:chemical-potential:10}
$x_1\,\dd(Ax_2^2) + x_2\,\dd(Ax_1^2) = 2Ax_1x_2(\dd x_2 + \dd x_1) = 0$, want
$x_1 + x_2 = 1$. Als $x_1 \to 0$, gaat $\gamma_1 \to \mathrm e^{A}$ en is $p_1 =
\gamma_1x_1p_1^* \approx \mathrm e^{A}p_1^*\,x_1$: $k_{H,1} = p_1^*\mathrm e^A$.
Voor $A = 1.1$ is $\mathrm e^{1.1} = 3.0$: de helling van Henry is driemaal die van Raoult.
\end{solution}

\begin{solution}{exo:b2:chemical-potential:11}
(a) Een vluchtige opgeloste stof draagt haar eigen dampdruk bij; de damp is niet
langer zuiver oplosmiddel en de afleiding (zuivere damp van het oplosmiddel in
evenwicht met de oplossing) gaat niet op. (b) Als de opgeloste stof in de vaste fase opgenomen wordt, daalt
ook de \omterm{def:b2:chemical-potential:chemical-potential}{chemische potentiaal} van de vaste stof, met $RT\ln x_1^{\text{s}}$. Als de
vaste oplossing dezelfde samenstelling heeft als de vloeistof, schuiven beide lijnen
evenveel omlaag en verandert het vriespunt niet: de daling
vereist een opgeloste stof die het kristal uitstoot.
\end{solution}

\begin{solution}{exo:b2:chemical-potential:12}
Cryoscopie: $b_{\min} = 0.01/1.86 = \qty{5.4e-3}{mol/kg}$; ebullioscopie:
$0.01/0.513 = \qty{1.9e-2}{mol/kg}$. Bij $\qty{5.4e-3}{mol/L} =
\qty{5.4}{mol/m^3}$ is $\Pi = 5.4 \times 8.314 \times 298 = \qty{1.3e4}{Pa}$,
meer dan een meter water; een osmometer leest enkele pascal af. Osmometrie is
verreweg het gevoeligst; daarom meet ze de grote molaire massa's van
polymeren, waarvan de molaliteiten minuscuul zijn.
\end{solution}

\begin{solution}{pb:b2:chemical-potential:1}
\textbf{1.} $\mu_1 = \mu_1^*(T, p) + RT\ln x_1$.
\textbf{2.} Diol $30/62.07 = \qty{0.483}{mol}$, water $70/18.015 =
\qty{3.886}{mol}$; $x_1 = 3.886/4.369 = 0.889$.
\textbf{3.} $p = 0.889 \times 3.17 = \qty{2.82}{kPa}$.
\textbf{4.} $b = 0.483/0.070 = \qty{6.90}{mol/kg}$.
\textbf{5.} Het diol kookt bij \qty{197}{\celsius}: zijn dampdruk bij
kamertemperatuur ligt ver onder die van water.
\textbf{6.} $\mu_1^{\text{s}}(T_f) = \mu_1^{*\text{l}}(T_f) + RT_f\ln x_1$.
\textbf{7.} $\ln x_1 = -\Delta_{\text{fus}}G(T)/(RT)$ met
$\Delta_{\text{fus}}G = \mu_1^{*\text{l}} - \mu_1^{\text{s}}$, nul bij $T_f^*$;
Gibbs--Helmholtz, $\dd(\Delta_{\text{fus}}G/T)/\dd T = -\Delta_{\text{fus}}H/T^2$,
geïntegreerd van $T_f^*$ tot $T_f$, geeft $\Delta_{\text{fus}}G(T_f)/T_f =
\Delta_{\text{fus}}H(1/T_f - 1/T_f^*)$, waaruit de betrekking volgt.
\textbf{8.} $\ln x_1 \approx -x_2 \approx -bM_1$ en $1/T_f - 1/T_f^* \approx
-\Delta T_f/T_f^{*2}$: $\Delta T_f = (RT_f^{*2}M_1/\Delta_{\text{fus}}H)\,b$.
\textbf{9.} $333.4 \times 18.015 = \qty{6007}{J/mol}$; $K_f = 8.314 \times
273.15^2 \times 0.018015/6007 = \qty{1.86}{K.kg/mol}$.
\textbf{10.} Dat $\Delta_{\text{fus}}H$ niet verandert tussen $T_f$ en
$T_f^*$ (dat doet ze wel, volgens Kirchhoff, hier ongeveer 2\,\% per \qty{3}{K}).
\textbf{11.} $\Delta T_f = 1.86 \times 6.90 = \qty{12.8}{K}$: \qty{-12.8}{\celsius}.
\textbf{12.} $1/T_f = 1/273.15 - (8.314/6007)\ln 0.889$: $T_f = \qty{261.6}{K}$,
\qty{-11.6}{\celsius}.
\textbf{13.} De exacte ideale betrekking, die geen verdunde oplossing
veronderstelt (hier is 11\,\% van de moleculen diol). Beide veronderstellen een \omterm{def:b2:chemical-potential:ideal-mixture}{ideaal mengsel}
en een constante smeltenthalpie; water en diol, verbonden door waterstofbruggen,
zijn niet ideaal, dus het echte vriespunt wijkt af.
\textbf{14.} Zuiver ijs: het diol blijft in de vloeistof.
\textbf{15.} Water dat als ijs verdwijnt, verlaagt $x_1$ in de vloeistof, en volgens de
betrekking daalt de evenwichtstemperatuur: de koelvloeistof bevriest over een gebied
van temperaturen, niet in één punt.
\textbf{16.} $K_b = 8.314 \times 373.12^2 \times 0.018015/40\,650 =
\qty{0.513}{K.kg/mol}$.
\textbf{17.} $\Delta T_b = 0.513 \times 6.90 = \qty{3.5}{K}$.
\textbf{18.} $1/T = 1/373.12 - (8.314/40\,650)\ln(2.0/1.013)$: $T = \qty{393.5}{K}$,
ongeveer \qty{120}{\celsius}.
\textbf{19.} De dop: \qty{20}{K} tegenover \qty{3.5}{K} voor het diol.
\textbf{20.} $x_1 = \exp[-(6007/8.314)(1/253.15 - 1/273.15)] = 0.811$.
\textbf{21.} $x_2 = 0.189$: $w = 0.189 \times 62.07/(0.189 \times 62.07 + 0.811
\times 18.015) = 0.44$.
\textbf{22.} Verdunde wet: $b = 20/1.86 = \qty{10.75}{mol/kg}$, $w = 10.75 \times
62.07/(1000 + 10.75 \times 62.07) = 0.40$. De linearisaties ($\ln x_1 \approx
-x_2$, $x_2 \approx bM_1$) overschatten de verlaging per mol diol bij zo'n
hoge concentratie, zodat de verdunde wet om minder diol vraagt.
\textbf{23.} In het model van het ideale mengsel blijft een koelvloeistof met een massafractie
diol $\boldsymbol{w \approx 0.44}$ vloeibaar tot \qty{-20}{\celsius}.
\end{solution}
